\section*{第 \ref{ch:g3:temperature-thermometers} 章 --- 温度与温度计}

\begin{solution}{exo:g3:temperature-thermometers:1}
同样的冷热，温度计给出的示数总是相同，不论是谁、
来读、何时读。手被三只碗骗了（同一温水
同时让人觉得暖和又觉得凉），还有金属
在同样的阴凉处比木头摸起来更凉。
\end{solution}

\begin{solution}{exo:g3:temperature-thermometers:2}
水在 \qty{0}{\celsius} 时凝固，在 \qty{100}{\celsius} 时沸腾，
均以摄氏度计。
\end{solution}

\begin{solution}{exo:g3:temperature-thermometers:3}
受热时，封闭的液体膨胀，液柱上升；
遇冷时，液体收缩，液柱下沉。玻璃管很细，
即使极小的膨胀也能使液柱明显上升。
\end{solution}

\begin{solution}{exo:g3:temperature-thermometers:4}
读数过早（液柱尚未停止移动），
以及斜着看，而不是让视线与液柱顶端齐平。
\end{solution}

\begin{solution}{exo:g3:temperature-thermometers:5}
沸水：\qty{100}{\celsius}。舒适的客厅：
\qty{20}{\celsius}。冰箱：\qty{4}{\celsius}。人体：
\qty{37}{\celsius}。
\end{solution}

\begin{solution}{exo:g3:temperature-thermometers:6}
锅更热，高出 $100 - 35 = 65$ 度，
烫得不能触摸：锅的温度正是沸腾的温度。
\end{solution}

\begin{solution}{exo:g3:temperature-thermometers:7}
零下三摄氏度，写作 \qty{-3}{\celsius}。它
比 \qty{-8}{\celsius} \emph{更暖和}：零下八度
比零下三度更深入冰冻。
\end{solution}

\begin{solution}{exo:g3:temperature-thermometers:8}
$14 - 6 = 8$：早晨气温升高了 $8$ 度。
\end{solution}

\begin{solution}{exo:g3:temperature-thermometers:9}
夜晚气温降到 \qty{0}{\celsius} 以下——水的凝固刻度。
结冰的水洼就是一支由天气做成的温度计。
\end{solution}

\begin{solution}{exo:g3:temperature-thermometers:10}
不对——\qty{16}{\celsius} 偏凉，不算舒适。艾玛的皮肤
是拿走廊和她刚离开的寒冷花园作比较，所以她觉得暖和。
家人应当相信温度计。温度计比较的是
刻度上固定的标记，
而不是它刚刚所在的地方。
\end{solution}

\begin{solution}{exo:g3:temperature-thermometers:11}
沸水无论沸腾多久，都保持在 \qty{100}{\celsius}——
这个刻度是一个上限：继续加热只会产生更多
水蒸气，而不会使水更热。这正是 \qty{100}{\celsius} 能被
选作可靠刻度的原因。
\end{solution}
